%!TeX root = ../incremental_SS_Translation.tex

\chapter{Śārīrasthāna 6:  Each Individual Vital Part of the Body}


\section{Introduction}
\texttt{Draft text by Jan Gerris.
    The Sanskrit text of this adhyāya is a lightly 
    edited transcript of NAK 5-333,
    waiting for full critical attention.}

\section{Literature} 

Meulenbeld offered an annotated overview of this chapter and a
bibliography of earlier scholarship to 2002 and, in his notes, citations
of the parallel passages in the \CS.\fvolcite{IA}[254--255]{meul-hist}   

The chapter on marmans in Vāgbhaṭa's \AH\ was translated by \citet{wuja-2003}.

The first draft of the following translation was prepared by Jan
Gerris.
% and later revised in collaboration Philipp Maas, John Nemec,
%Dominik Wujastyk, Madhusudan Rimal, Kenneth Zysk and Andrey Klebanov.

\newpage

\section{Translation}

\begin{translation}

    % First draft by Jan Gerris, [date]
  
\bigskip       

 \begin{tt}
    \raggedright
    \bigskip
    \marginpar{Draft tr.\ from here}

           

\item[1]
Now, therefore, we shall explain the anatomical description of each individual vital part of the body. 

\item[2] missing
\item[3]
There are one hundred and seven vital points.
These vital points are of five types.
They are as follows:
Muscle-vital points, vessel-vital points, ligament-vital points, bone-vital points, and joint-vital points.  Indeed, there are no other vital points apart from muscle, vessels, ligaments, bones, and joints. Because they are not found. 

[cf. 1938 ed. 3.6.4
Among them, there are eleven muscle-marmas and forty-one vessel-marmas. 
There are twenty-seven ligament-marmas.
There are eight bone-marmas. 
There are twenty joint-marmas.
This is that one hundred and seven vital points.

\item[5]
Of these, eleven are located in one leg. 
By this, the other leg and the two arms are also explained. 
There are twelve in the abdomen and chest.
Fourteen in the back.
There are thirty-seven in the neck and above it.

\item[6]
The vital points of the leg are as follows: Kṣipra, Talahṛdaya, Kūrca, Kūrcaśiras, Gulpha, Indravasti, Jānu, Āṇi, Ūrvī, Lohitākṣa, and Viṭapa.  By this, the other leg and the two arms are explained.  Next, we shall describe the vital points of the abdomen and the thorax. They are the rectum, the bladder, the navel, the heart, the two Stanamūlas, the two Stanarohitas, the two Apalāpas, and the two Apastambhas.  The vital points of the back are as follows: Kaṭīkataruṇa, Kukundara, Nitamba, Pārśvasandhi, Bṛhatī, Aṃsaphalaka, and Aṃsa. Above the collarbone, there are four Dhamanīs, eight Mātrikās, two Kṛkāṭikās, two Vidhuras, two Phaṇas, two Apāṅgas, two Śaṅkhas, two Āvartas, two Utkṣepas, one Sthapanī, five Sīmantas, four Śṛṅgāṭakas, and one Adhipati.

\item[7]
Among those, Talahṛdaya, Indravasti, Guda, and Stanarohita are the muscle marmas.
Nīladhamanī, Mātrikā vessels, Śṛṅgāṭaka, Apāṅga, Sthapanī, Phaṇa, Stanamūla, Apalāpa, Apastambha, Hṛdaya, Nābhi, Pārśvasandhi, Bṛhatī, Lohitākṣa, and Ūrvī are the vascular marmas.
Viṭapa, Kakṣadhara, Kūrca, Kūrcaśiras, Vasti, Kukṣi, Aṃsa, Vidhura, and Utkṣepa are the ligament marmas. 
Kaṭīkataruṇa, Nitamba, Aṃsaphalaka, and Śaṅkha are the bone marmas.
Jānu, Kūrpara, Sīmanta, Adhipati, Gulpha, Maṇibandha, Kukundara, Āvarta, and Kṛkāṭikā are the joint marmas.

\item[8]
These vital points are of five distinct types.
They are as follows:
Those that cause instantaneous death
those that cause death after a lapse of time
those that kill upon the extraction of a foreign body
those that cause disability
and those that cause pain. 
Among them, nineteen are of the type that cause instantaneous death. 
Thirty-three cause death after a lapse of time. 
Three are fatal upon extraction of a foreign body. 
Forty-four cause disability. 
Eight are said to cause pain. 


\item[9]
And there are verses on this subject.
The Śṛṅgāṭakas, the Adhipati, the two Śaṅkhas, the Kanṭhasirās, the Guda, the Hṛdaya, and the Vasti, along with the Nābhi, kill instantly if they are injured. 

\item[10]
The marmas of the chest, the Sīmantas, the Tala, Kṣipra, Indravasti, the Kaṭīkataruṇas, the Sandhis, the Pārśvaja, and the Bṛhatī,

\item[11]
and the two Nitambas—these are the ones that take life after a lapse of time. One should point out the two Utkṣepas and the Sthapanī as those that kill upon extraction of a foreign body.

\item[12]
The Lohitākṣa, Jānu, Urvī, Kūrca, Viṭapa, and Kūrpara, the Kukundare, Kakṣadhare, Vidura, and along with the Kṛkātike


\item[13]
The two aṃsas, the two aṃsaphalakas, the two apāṅgas, the two nīlas, the two manyās, and the two phaṇas are likewise called disability-causing marmas, as are the two āvartas.


\item[14]
A wise person should know that these eight—the two gulphas, the two maṇibandhas, and the four kūrcaśiras—are pain-causing marmas.
Five distinct classifications of marmas have been stated according to their effects.


\item[15]
Marmas are indeed the meeting points of muscle, vessels, ligaments, bones, and joints.
Life-breaths reside in them by their very nature.  When injured in the marmas, they undergo various states.


\item[16]
Among those, there are the sadyah-prāṇahara (instantly fatal) marmas.
They are fiery in nature.
When the fiery qualities are exhausted, they destroy life.
The viśalyaghna (fatal upon extraction) marmas are airy in nature; as long as the air remains blocked by the head of the weapon, the person lives.
However, as soon as the weapon is extracted, the air residing in the vital spot escapes.
Therefore, one lives while the weapon is in place, but dies once the weapon is removed.
The disability-causing marmas are lunar in nature.
Because of its coolness and stability, Soma (the moon) provides support for the life-breaths. 
The pain-causing marmas are composed mostly of fire and air; specifically, those two cause pain, though others say pain is composed of all five elements.


\item[17]
Some claim that the vital points which cause immediate death result from the combination of all five tissues—flesh and the rest—in their fully developed state. 
Those lacking one tissue or where they are present in small measure cause death after a lapse of time through their combination. 
However, those lacking two tissues cause death upon the extraction of a foreign object.
Those lacking three tissues cause functional deformity. 
When only one tissue is involved, they cause only pain.
Because of this, bleeding occurs even when the vital points of the bones are injured.


\item[18]
And here are the verses on this subject:
The four types of vessels in the body are generally situated within the vital points; they nourish the ligaments, bones, flesh, and joints, thereby sustaining the body.


cf. 1938 ed. 3.6.19]
Then, when a vital point is wounded, the vital air (Vāyu) increases and spreads in all directions. 
As that all-pervading air expands, it spreads intense pains throughout the body. 

\item[20]
Overwhelmed by pain, the body then collapses and the person's consciousness perishes.

\item[21]
By this, the remainder is explained.

\item[22]
Among these, [one type] takes life instantly; [another], when struck/bitten at the extremity, kills after an interval of time. A vital point that causes death after a lapse of time, if pierced at its edge, results in deformity. The "Vishalyaghna" type of vital point also results in deformity. They produce various types of intense pain. A deformity-causing vital point causes distress after a lapse of time. Or it produces an ache, and a pain-causing point produces mild pain.

[cf. 1938 ed. 3.6.23
Among those, the vital points that cause immediate death kill within seven nights. 
The vital points that cause death after a lapse of time kill after a fortnight or a month.
Even among those, the "Kshipra" vital points may sometimes kill quickly. 
The "Vishalyaghna" and "Vaikalyakara" vital points may also cause death if they are excessively injured.

\item[24]

Hereafter, we shall describe the vital points of the lower limb.
Among those, the vital point named Kṣipra is located between the big toe and the second toe of the foot.
Following the middle toe, in the center of the sole of the foot, is the vital point named Talahṛdaya; pain there results in death. 
One finger-breadth above the Kṣipra, on both sides, is the vital point named Kūrca.
Below the ankle joint is the vital point named Kūrcaśiras; pain and swelling occur there.
At the junction of the foot and the calf is the vital point named Gulpha; an injury there causes pain, stiffness of the leg, or lameness.
Opposite the heel, in the middle of the calf, is the vital point named Indrabasti; death occurs there due to loss of blood.
At the junction of the calf and the thigh is the vital point named Jānu; lameness occurs there.
Three finger-breadths above the knee is the vital point named Āṇi; there, an increase in swelling and stiffness of the leg occurs.
In the middle of the thigh is the vital point named Urvī; due to loss of blood there, wasting of the leg occurs.
Above the Urvī and below the groin, at the root of the thigh, is the vital point named Lohitākṣa; death or paralysis occurs there due to loss of blood.
Between the groin and the scrotum is the vital point named Viṭapa; impotence and deficiency of semen occur there.
Thus, these eleven vital points of the leg have been explained.
By this, the other leg and the two arms are also explained. 
Specifically, those which are the Gulpha, Jānu, Vaṅkṣaṇa, and Viṭapa in the leg are the Maṇibandha, Kūrpara, and Kakṣadhara in the arm.
Just as the vital point named Gulpha is at the junction of the foot and calf,
so too is the vital point named Maṇibandha at the junction of the palm and the forearm.
Just as the vital point named Jānu is at the junction of the calf and the thigh,
so too is the vital point named Kūrpara at the junction of the forearm and the upper arm.
Just as the vital point named Viṭapa is between the groin and the scrotum,
so too is the vital point named Kakṣadhara between the clavicle and the armpit
When these are pierced, the same complications occur.
Specifically, there is stiffness in the Maṇibandha,
paralysis in the Kūrand lameness in the Kakṣadhara.
Thus, these forty-four vital points of the extremities have been explained.

\item[25]
Hereafter, we shall describe the twelve vital points of the abdomen and the thorax.Among them, the one named Mahāguda is attached to the large intestine and serves as the expeller of wind and faeces. In that case, death is immediate. It contains little flesh and blood and is located internally. The vital point named Basti is the urinary bladder located in the pelvic region. There too, death is immediate.
Except in the case of a wound from a urinary stone, death occurs if it is pierced on both sides; if it is ruptured on one side, a wound discharging urine occurs.
However, if treated with effort, it heals.
The vital point named Nābhi is the origin of the vessels located between the stomach and the large intestine; there too, death is immediate.
The vital point named Hṛdaya is located in the thorax between the two breasts at the opening of the stomach, and is the seat of Sattva, Rajas, and Tamas; there too, death is immediate.
The two vital points named Stanamūla are located two finger-breadths below each of the breasts.
In that case, death occurs due to the thorax being filled with Kapha.
The two vital points named Stanarohita are above the nipples of the breasts.
In that case, death occurs due to the thorax being filled with blood, accompanied by cough and asthma.
The two vital points named Apalāpa are located below the shoulders and above the sides.
In that case, death occurs when the blood attains a state of suppuration.
On both sides of the thorax are the air-carrying channels.
These are the two vital points named Apastambha.
In that case, death occurs due to the thorax being filled with wind, accompanied by cough and asthma.
Thus, these twelve vital points of the abdomen and thorax have been explained. 

\item[26]
Hereafter, we shall describe the vital points of the back. 
There, on both sides of the spinal column, in the bones of the pelvic region, are the two vital points named Kaṭīkataruṇa.
In that case, death occurs from loss of blood, or the patient becomes pale, discoloured, and wasted.
On the outer part of the hips, on both sides of the spinal column, are the two vital points named Kukundara.
In that case, there is loss of sensation and weakness in the lower body, and loss of movement.
Above the pelvic bones, covering the stomach and attached to the inner sides, are the two vital points named Nitamba.
In that case, there is wasting of the lower body, weakness, and death.
Attached to the lower inner sides, in the middle of the hip and side, diagonally above the hip, are the two vital points named Pārśvasandhi.
In that case, the patient dies due to the thorax being filled with blood.
On both sides of the spinal column, in a straight line from the roots of the breasts, are the two vital points named Bṛhatī.
On the back, on both sides of the vertebral column.
At the connection of the trika (sacrum) are the two vital spots named aṃsaphalaka.
In that instance, there is wasting and loss of movement in the arms.
Between the head of the arm, the neck, and the shoulder-blade, at the junction of the shoulder, are the two vital spots named kṛkāṭikā.
In that case, the arm becomes stiff or the back becomes broken.
Thus, these fourteen vital spots of the back have been explained.
\item[27]

Hereafter, we shall describe the vital points located in the region above the collarbones.
There, on both sides of the windpipe, are four channels arranged crosswise.
Two are named Nīlā and two are named Manyā. 
In that case, muteness, alteration of the voice, and loss of the sense of taste occur. 
In the neck, on both sides, are four vital points called Sirāmātṛkā.
In that case, death is immediate.
At the junction of the head and neck are two vital points called Kṛkāṭikā.
In that case, there is a shaking of the head or death.
Located below and behind the ears are the two vital points named Vidhura.
In that case, deafness occurs. 
On both sides of the nasal passage, connected to the internal channels, are the two vital points named Phaṇā. 
In that case, it causes the loss of the sense of smell. 
Below the ends of the tails of the eyebrows and outside the eyes are the two vital points named Apāṅga.
In that case, there is blindness or impairment of vision.
Above the tails of the eyebrows, between the ears and the forehead, are the two vital points named Śaṅkha.
In that case, death is immediate. 
Above the eyebrows, in the middle of the depressions, are the two vital points named Āvarta.
There also, blindness or impairment of vision occurs. 
Above the Āvarta and Śaṅkha points, at the hairline, are the two vital points named Utkṣepa.
In that case, one lives as long as the shaft remains, or if the shaft falls out due to suppuration, but not if it is extracted. 
Between the eyebrows is the vital point named Sthapanī.
Like the Utkṣepa, there are five joints divided on the head called Sīmanta.
In that case, death occurs due to madness, fear, and loss of consciousness. 
In the midst of the vessels that nourish the nose, ears, eyes, and tongue, is the meeting point of the head's channels; these are the four Śṛṅgāṭaka vital points.
There also, death is immediate.
Inside the head, at the top, where the vessels and joints meet at the hair-whorl, is the vital point named Adhipati. 
There also, death is immediate.
Thus, these thirty-seven vital points situated in the region above the collarbones have been explained.


\item[28]
And there are verses on this subject. 
The Urvī, Śiras, Viṭapa, along with the Kakṣa and Pārśva, are each measured as one finger-breadth; the Stanamūla and Gulpha are known to be two finger-breadths; the Maṇibandha is the sixth; and three fingers are the measure for the Jānu along with the Kūrpara.

\item[29]
They say there are [vital spots] in the heart, bladder, kūrca, anus, and navel, as well as four in the head, five in the throat, and twelve [others] These [measure] as much as one's own fist when contracted; know the remaining ones to be the size of a thumb-joint or half a finger.


\item[31]
When the hands or feet are severed, the blood of men flows only in small amounts because of the contraction of the openings of the vessels. Therefore, when hands and feet are severed, they [the vessels] become like lotuses in stalks that have been crushed in a circular motion.

\item[32]
In the case of [injuries to] the kṣipra and talsahṛdaya vital spots, blood flows out due to the defect of the piercing, and the vāyu causes great pain. Indeed, those pierced there perish just as lotuses do from the crushing of their filaments and petals. 


\item[33]
The vital spots are said to be half the scope of surgery, because mortals struck in the vital spots do not [usually] survive. 

\item[34]
If some survive there by the skill of the physician, they undoubtedly suffer from deformity.
People survive even with shattered or crushed abdominal cavities or skull bones, and with body parts struck by weapons.


\item[35]
Those whose thighs, arms, feet, or hands are completely severed survive if the various blows have not fallen upon the vital spots.
Fire, Moon, Wind, and the qualities of Sattva, Rajas, and Tamas are generally present in the vital spots of men, along with the individual soul.

\item[36]
Therefore, living beings do not survive when struck in the vital points; there is a lack of perception regarding sense objects, and a derangement of the mind and intellect. 


\item[37]
Intense and various pains occur when a person is struck in a way that takes away their life. 
Even when struck in a point that causes death after a lapse of time, the depletion of the bodily constituents is certain for men. 

\item[38]
Then, due to that depletion of the elements and because of the agonies, the creature perishes. When a point that causes deformity is struck, only through the skill of a physician...

\item[39]
...having incurred the wasting of the body, one would attain a state of deformity. Regarding the points that are fatal when a foreign object is extracted, the reason previously mentioned should be understood.

\item[40]
The pain-causing vital points, when wounded, produce various pains and eventually lead to deformity if the patient falls under the control of a bad physician.

\item[41]
One should understand that an injury to the vital points resulting from excision, incision, impact, burning, or even tearing, exhibits identical symptoms.
\item[42]
Is there any injury to a vital spot that is of little consequence or entirely without danger? Generally, those struck in the vital spots either suffer permanent deformity or meet with death.

\item[43]
Indeed, the various disorders that take root in the vital spots and spread through the human body are usually the most difficult to cure, even when treated with great effort.

Thus ends the sixth chapter in the Section on Anatomy.

\end{tt}
\end{translation}
